% ---------------------------------------------------------------------------
% Study report · Do three thermometers tell the same story?
%
% Build:  pdflatex proxy_comparison_report.tex  (twice, for the ToC)
% Figures are read from ../outputs/ and are produced by the paired notebook.
% ---------------------------------------------------------------------------
\documentclass[11pt,a4paper]{article}

\newcommand{\runninghead}{Proxy comparison · on-structure, ground station, ERA5}

% ---------------------------------------------------------------------------
% reportstyle.tex — shared preamble for the study reports under studies/.
%
% Kept deliberately inside the package set shipped with TeX Live "basic":
% no tcolorbox, no siunitx, no enumitem, no titlesec. Everything below
% compiles with a stock pdflatex installation.
%
% Usage, from a report directory one level below a study folder:
%     % ---------------------------------------------------------------------------
% reportstyle.tex — shared preamble for the study reports under studies/.
%
% Kept deliberately inside the package set shipped with TeX Live "basic":
% no tcolorbox, no siunitx, no enumitem, no titlesec. Everything below
% compiles with a stock pdflatex installation.
%
% Usage, from a report directory one level below a study folder:
%     % ---------------------------------------------------------------------------
% reportstyle.tex — shared preamble for the study reports under studies/.
%
% Kept deliberately inside the package set shipped with TeX Live "basic":
% no tcolorbox, no siunitx, no enumitem, no titlesec. Everything below
% compiles with a stock pdflatex installation.
%
% Usage, from a report directory one level below a study folder:
%     \input{../../_shared/reportstyle.tex}
%     \graphicspath{{../outputs/}}
% ---------------------------------------------------------------------------

\usepackage[T1]{fontenc}
\usepackage[utf8]{inputenc}
\usepackage{textcomp}
\usepackage{lmodern}
\usepackage{microtype}
\usepackage[margin=2.4cm,headheight=15pt]{geometry}
\usepackage{graphicx}
\usepackage{booktabs}
\usepackage{tabularx}
\usepackage{array}
\usepackage{amsmath}
\usepackage{xcolor}
\usepackage[font=small,labelfont=bf,justification=justified]{caption}
\usepackage{float}
\usepackage{fancyhdr}
\usepackage[hidelinks]{hyperref}

% --- Palette ---------------------------------------------------------------
\definecolor{rulegrey}{HTML}{B9C0C7}
\definecolor{fillgrey}{HTML}{F2F4F6}
\definecolor{failred}{HTML}{A5202B}
\definecolor{passgreen}{HTML}{1B6B3A}
\definecolor{accent}{HTML}{1F4E79}

% --- Semantic shorthands ---------------------------------------------------
\newcommand{\degC}{\textdegree C}
\newcommand{\mdegC}{mdeg/\textdegree C}
\newcommand{\fail}{\textcolor{failred}{\textbf{fail}}}
\newcommand{\pass}{\textcolor{passgreen}{\textbf{pass}}}
\newcommand{\worse}{\textcolor{failred}{worse}}
\newcommand{\better}{\textcolor{passgreen}{better}}

% --- Highlighted panel -----------------------------------------------------
% #1 = heading, #2 = body
\newcommand{\panel}[2]{%
  \begin{center}%
  \setlength{\fboxsep}{9pt}%
  \fcolorbox{rulegrey}{fillgrey}{%
    \begin{minipage}{0.93\textwidth}%
      {\bfseries\color{accent}#1}\par\medskip
      #2
    \end{minipage}}%
  \end{center}}

% --- Numeric column: right aligned, fixed width ----------------------------
\newcolumntype{R}[1]{>{\raggedleft\arraybackslash}p{#1}}
\newcolumntype{L}[1]{>{\raggedright\arraybackslash}p{#1}}
\newcolumntype{Y}{>{\raggedright\arraybackslash}X}

% --- Table look ------------------------------------------------------------
\renewcommand{\arraystretch}{1.15}
\setlength{\tabcolsep}{6pt}

% --- Headers ---------------------------------------------------------------
\pagestyle{fancy}
\fancyhf{}
\fancyhead[L]{\small\itshape\runninghead}
\fancyhead[R]{\small\thepage}
\renewcommand{\headrulewidth}{0.4pt}

% --- Section spacing without titlesec --------------------------------------
\makeatletter
\renewcommand\section{\@startsection{section}{1}{\z@}%
  {-3.2ex \@plus -1ex \@minus -.2ex}{2.0ex \@plus.2ex}%
  {\normalfont\Large\bfseries\color{accent}}}
\renewcommand\subsection{\@startsection{subsection}{2}{\z@}%
  {-2.6ex \@plus -1ex \@minus -.2ex}{1.4ex \@plus.2ex}%
  {\normalfont\large\bfseries}}
\makeatother

% --- Figure defaults -------------------------------------------------------
\setkeys{Gin}{width=\linewidth}
\captionsetup[figure]{skip=6pt}

%     \graphicspath{{../outputs/}}
% ---------------------------------------------------------------------------

\usepackage[T1]{fontenc}
\usepackage[utf8]{inputenc}
\usepackage{textcomp}
\usepackage{lmodern}
\usepackage{microtype}
\usepackage[margin=2.4cm,headheight=15pt]{geometry}
\usepackage{graphicx}
\usepackage{booktabs}
\usepackage{tabularx}
\usepackage{array}
\usepackage{amsmath}
\usepackage{xcolor}
\usepackage[font=small,labelfont=bf,justification=justified]{caption}
\usepackage{float}
\usepackage{fancyhdr}
\usepackage[hidelinks]{hyperref}

% --- Palette ---------------------------------------------------------------
\definecolor{rulegrey}{HTML}{B9C0C7}
\definecolor{fillgrey}{HTML}{F2F4F6}
\definecolor{failred}{HTML}{A5202B}
\definecolor{passgreen}{HTML}{1B6B3A}
\definecolor{accent}{HTML}{1F4E79}

% --- Semantic shorthands ---------------------------------------------------
\newcommand{\degC}{\textdegree C}
\newcommand{\mdegC}{mdeg/\textdegree C}
\newcommand{\fail}{\textcolor{failred}{\textbf{fail}}}
\newcommand{\pass}{\textcolor{passgreen}{\textbf{pass}}}
\newcommand{\worse}{\textcolor{failred}{worse}}
\newcommand{\better}{\textcolor{passgreen}{better}}

% --- Highlighted panel -----------------------------------------------------
% #1 = heading, #2 = body
\newcommand{\panel}[2]{%
  \begin{center}%
  \setlength{\fboxsep}{9pt}%
  \fcolorbox{rulegrey}{fillgrey}{%
    \begin{minipage}{0.93\textwidth}%
      {\bfseries\color{accent}#1}\par\medskip
      #2
    \end{minipage}}%
  \end{center}}

% --- Numeric column: right aligned, fixed width ----------------------------
\newcolumntype{R}[1]{>{\raggedleft\arraybackslash}p{#1}}
\newcolumntype{L}[1]{>{\raggedright\arraybackslash}p{#1}}
\newcolumntype{Y}{>{\raggedright\arraybackslash}X}

% --- Table look ------------------------------------------------------------
\renewcommand{\arraystretch}{1.15}
\setlength{\tabcolsep}{6pt}

% --- Headers ---------------------------------------------------------------
\pagestyle{fancy}
\fancyhf{}
\fancyhead[L]{\small\itshape\runninghead}
\fancyhead[R]{\small\thepage}
\renewcommand{\headrulewidth}{0.4pt}

% --- Section spacing without titlesec --------------------------------------
\makeatletter
\renewcommand\section{\@startsection{section}{1}{\z@}%
  {-3.2ex \@plus -1ex \@minus -.2ex}{2.0ex \@plus.2ex}%
  {\normalfont\Large\bfseries\color{accent}}}
\renewcommand\subsection{\@startsection{subsection}{2}{\z@}%
  {-2.6ex \@plus -1ex \@minus -.2ex}{1.4ex \@plus.2ex}%
  {\normalfont\large\bfseries}}
\makeatother

% --- Figure defaults -------------------------------------------------------
\setkeys{Gin}{width=\linewidth}
\captionsetup[figure]{skip=6pt}

%     \graphicspath{{../outputs/}}
% ---------------------------------------------------------------------------

\usepackage[T1]{fontenc}
\usepackage[utf8]{inputenc}
\usepackage{textcomp}
\usepackage{lmodern}
\usepackage{microtype}
\usepackage[margin=2.4cm,headheight=15pt]{geometry}
\usepackage{graphicx}
\usepackage{booktabs}
\usepackage{tabularx}
\usepackage{array}
\usepackage{amsmath}
\usepackage{xcolor}
\usepackage[font=small,labelfont=bf,justification=justified]{caption}
\usepackage{float}
\usepackage{fancyhdr}
\usepackage[hidelinks]{hyperref}

% --- Palette ---------------------------------------------------------------
\definecolor{rulegrey}{HTML}{B9C0C7}
\definecolor{fillgrey}{HTML}{F2F4F6}
\definecolor{failred}{HTML}{A5202B}
\definecolor{passgreen}{HTML}{1B6B3A}
\definecolor{accent}{HTML}{1F4E79}

% --- Semantic shorthands ---------------------------------------------------
\newcommand{\degC}{\textdegree C}
\newcommand{\mdegC}{mdeg/\textdegree C}
\newcommand{\fail}{\textcolor{failred}{\textbf{fail}}}
\newcommand{\pass}{\textcolor{passgreen}{\textbf{pass}}}
\newcommand{\worse}{\textcolor{failred}{worse}}
\newcommand{\better}{\textcolor{passgreen}{better}}

% --- Highlighted panel -----------------------------------------------------
% #1 = heading, #2 = body
\newcommand{\panel}[2]{%
  \begin{center}%
  \setlength{\fboxsep}{9pt}%
  \fcolorbox{rulegrey}{fillgrey}{%
    \begin{minipage}{0.93\textwidth}%
      {\bfseries\color{accent}#1}\par\medskip
      #2
    \end{minipage}}%
  \end{center}}

% --- Numeric column: right aligned, fixed width ----------------------------
\newcolumntype{R}[1]{>{\raggedleft\arraybackslash}p{#1}}
\newcolumntype{L}[1]{>{\raggedright\arraybackslash}p{#1}}
\newcolumntype{Y}{>{\raggedright\arraybackslash}X}

% --- Table look ------------------------------------------------------------
\renewcommand{\arraystretch}{1.15}
\setlength{\tabcolsep}{6pt}

% --- Headers ---------------------------------------------------------------
\pagestyle{fancy}
\fancyhf{}
\fancyhead[L]{\small\itshape\runninghead}
\fancyhead[R]{\small\thepage}
\renewcommand{\headrulewidth}{0.4pt}

% --- Section spacing without titlesec --------------------------------------
\makeatletter
\renewcommand\section{\@startsection{section}{1}{\z@}%
  {-3.2ex \@plus -1ex \@minus -.2ex}{2.0ex \@plus.2ex}%
  {\normalfont\Large\bfseries\color{accent}}}
\renewcommand\subsection{\@startsection{subsection}{2}{\z@}%
  {-2.6ex \@plus -1ex \@minus -.2ex}{1.4ex \@plus.2ex}%
  {\normalfont\large\bfseries}}
\makeatother

% --- Figure defaults -------------------------------------------------------
\setkeys{Gin}{width=\linewidth}
\captionsetup[figure]{skip=6pt}

\graphicspath{{../outputs/}}

\title{\vspace{-1.2cm}\textbf{Do three thermometers tell the same story?}\\[2mm]
\large On-structure logger, town weather station and ERA5 reanalysis,
26 July 2018 to 13 August 2026}
\author{Study report · \texttt{studies/proxy\_comparison/}}
\date{13 August 2026}

\begin{document}
\maketitle
\thispagestyle{fancy}

\begin{abstract}
\noindent
Three independent sources now measure the environment at Gubbio and the project
has used only one of them. They share channel names and they do not measure the
same quantity: a housing on a sun-exposed wall, a standard screen in town, and an
average over roughly nine kilometres of Apennine terrain.

This comparison had to come before any of them is used to fill a gap. The
study \texttt{unified\_dataset} established that the on-structure channels
vanish with the target --- on 100\,\% of the hours the inclinometer is missing,
so are its air temperature, humidity and battery --- so an independent source is
the only escape. Whether that escape is real depends on whether the proxy carries
the same relationship with the wall.

\textbf{It largely does, and the availability is decisive.} During the hours the
inclinometer is missing, the ground station is present on \textbf{94.3\,\%} and
ERA5 on \textbf{99.6\,\%}. Against the on-structure sensor the ground station
agrees to a mean absolute \textbf{1.12\,\degC} and ERA5 to \textbf{1.53\,\degC},
with correlations of 0.988 and 0.978.

\textbf{One expectation was overturned.} A gridded product was expected to
under-represent the daily temperature swing. It does the opposite: ERA5's diurnal
amplitude is \textbf{19\,\% larger} than the on-structure sensor's and the ground
station's is 15\,\% larger. The sensor closest to the wall is the most
\emph{damped} of the three, which is what a housing with thermal mass on a
massive masonry wall should produce.

\textbf{Agreement collapses under differencing.} Correlation between sources
falls from 0.99 in levels to \textbf{0.70} in first differences for temperature,
and from 0.88 to 0.55 for humidity. The level agreement is largely a shared
annual cycle; the hour-to-hour movement, which is what the wall responds to,
agrees far less well.

\textbf{A defect is confirmed by triangulation.} The on-structure radiation
channel agrees with neither proxy ($r \approx 0.74$), while the two proxies agree
with each other at $r = 0.930$ and a slope of 0.998. Two independent instruments
agreeing with each other and disagreeing with a third identify which one is
wrong.

\textbf{The relationship with the wall had to be measured inside contiguous
blocks, and one defect had to be registered before it could be read at all.}
Computed on the whole record at once, as an earlier version of this report did,
the differenced relationship appears negligible for every source and the
on-structure radiation channel appears to \emph{reverse} sign. Both are
artefacts. Cut into blocks, kept inside one instrument era, and with 175 hours
of single-hour acquisition faults registered and masked --- hours in which the
inclination, the air temperature and the battery all fail together and recover
on the next sample --- air temperature reaches $r = -0.95$ in levels and
$-0.84$ in first differences against the calibrated inclination, and the
radiation reversal disappears.

\textbf{All three sources give the same sign}, no channel leads the deformation,
and every driver optimises at \textbf{zero transport delay}. Only solar
radiation needs an operator: a one- to two-hour thermal inertia lifts its
explained variance from 0.570 to \textbf{0.708}. Radiation adds one to three
points of explained variance beyond air temperature --- real, and small.
\end{abstract}

\tableofcontents

% ===========================================================================
\section{What is being compared, and why}
\label{sec:intro}

\subsection{Three sources, one name}

\begin{center}
\small
\begin{tabularx}{\textwidth}{l l Y}
\toprule
\textbf{Suffix} & \textbf{Source} & \textbf{What it measures} \\
\midrule
\texttt{\_str} & on-structure logger & Air inside the instrument housing on the
monitored wall \\
\texttt{\_gs} & ground station & Standard-exposure station in Gubbio town \\
\texttt{\_era5} & ERA5 reanalysis & Average over a grid cell of roughly nine
kilometres \\
\bottomrule
\end{tabularx}
\end{center}

\textbf{They are not replicates.} Disagreement between them is not error to be
minimised; it is the physical difference between three measurement situations,
and the study's job is to measure it and say what it costs.

\subsection{Why this had to come first}

\panel{The deadlock this study exists to break}{%
\texttt{studies/unified\_dataset/} established that every channel of the
monitored station fails at once. On \textbf{100\,\%} of the hours the
inclinometer is missing, the air temperature, humidity and battery are missing
too. The missingness is outage-driven, so no method conditioned on an
on-structure covariate can fill a real gap --- the covariates are gone with the
target.

\medskip
An independent source does not share the station's power supply. That is the only
escape available, and whether it works depends on two things: is the proxy
\emph{there} when the target is not, and does it carry the \emph{same
relationship} with the wall. This study answers both.}

\subsection{Two rules}

\textbf{The target is fixed.} The compensated inclination \texttt{inc\_comp} is
the calibrated reading: the manufacturer's formula applied with the temperature
measured at the transducer. That is an instrument calibration, not a regression.
The transducer experiences the temperature at the transducer, so substituting a
grid-cell or in-town temperature into the formula would be meaningless. Nothing
in this study re-estimates the coefficient and no alternative target is built.
That the on-structure temperature appears inside the target is a property of the
calibrated measurement, stated here and not designed around.

\textbf{The clock is measured, not assumed.} The sources do not share a time
convention, and an uncorrected offset would misplace every diurnal phase in the
study.

% ===========================================================================
\section{Harmonisation, and the clock}
\label{sec:clock}

\subsection{What arrived}

All three sources are placed on the hourly grid under one naming scheme,
\texttt{\{quantity\}\_\{source\}}, documented in
\texttt{docs/proxy-data-dictionary.md} --- which is generated from the mapping
tables in \texttt{pc\_lib} so that code and documentation cannot drift apart.

Two unit hazards were found and corrected. ERA5 reports relative humidity as a
\textbf{fraction} where the other two report per cent; left uncorrected it would
have appeared as a 99\,\% dry bias. The ground station is natively half-hourly
with 124 duplicated timestamps.

Values outside a plausible physical range are \textbf{reported, not removed}. A
radiation channel reading sixty watts at midnight is a finding about the sensor,
and deleting it would hide the finding.

\subsection{The radiation channels as measured}
\label{sec:series}

Every other figure in this report shows a statistic --- a mean profile, an
agreement scatter, a correlation. Figure~\ref{fig:series} shows the
measurements, and it is placed first because the strongest evidence against the
on-structure pyranometer is visible in the raw channel before any statistic is
computed.

The panels are restricted to the period the on-structure pyranometer covers,
21 February 2025 to 12 August 2026, because that is the only window over which
the three sources can be compared at all. The proxies run from 2018; that
difference in extent is what shapes the tiers in Section~\ref{sec:tiers} and is
reported there, and drawing it here would leave seven eighths of the
on-structure panel empty and hide the comparison this figure exists to make.

\begin{figure}[H]
\includegraphics{P_F09_radiation_series}
\caption{Solar radiation as measured, by source, over the on-structure record.
The upper three panels give one source each at hourly resolution; the black line
is a thirty-day rolling maximum, drawn as a seasonal ceiling rather than as a
smoothing of the values beneath it. Hours with no measurement are left as breaks
rather than joined, so an outage appears as the gap it is. The shaded band marks
the fortnight enlarged below, chosen as the fourteen days in which all three
sources are jointly most complete. The bottom panel overlays all three at full
hourly resolution over that fortnight, which is the only scale at which an
individual hourly value can be read.}
\label{fig:series}
\end{figure}

\textbf{The two proxies behave as a solar record should and the on-structure
channel does not.} The ground station and ERA5 panels each show one clean annual
arc, a ceiling near 1000\,W/m\textsuperscript{2} in summer falling to roughly
450 in January and recovering, on 99.2\,\% and 99.5\,\% of the window
respectively. The on-structure panel is present on \textbf{69.1\,\%} of the same
window and does none of that.

\textbf{It over-reads whenever it works.} Its peak is
\textbf{1358\,W/m\textsuperscript{2}} against 1004 for the ground station and
939 for ERA5, and in the enlarged fortnight its daily maxima run about
\textbf{1.5 times} both proxies, day after day, in clear and overcast conditions
alike. A clear-sky maximum above 1300\,W/m\textsuperscript{2} is not physically
available at this latitude.

\panel{The channel fails outright for two months, and under-reads for four more}{%
The monthly maxima make the failure explicit. Through the first summer the
on-structure channel runs above both proxies, as it does again in the second:
1291, 1317 and 1326\,W/m\textsuperscript{2} in May, June and July 2025 against
roughly 950 for each proxy, and 1353, 1312 and 1154 in June, July and August
2026.

\medskip
\textbf{In August and September 2025 it collapses.} Its monthly maxima are
\textbf{116} and \textbf{150}\,W/m\textsuperscript{2} while both proxies report
a normal late summer near \textbf{890}. A pyranometer that never exceeds
150\,W/m\textsuperscript{2} across an entire Umbrian August is not measuring
sunlight.

\medskip
\textbf{From December 2025 to March 2026 it under-reads}, inverting its usual
bias: 541, 496, 462 and 354\,W/m\textsuperscript{2} against proxy maxima that
rise from about 480 back to 816 as the season turns. By February and March it
sits \emph{below} both proxies, having been half again above them six months
earlier.

\medskip
\textbf{This matters for the agreement statistics of
Section~\ref{sec:absolute}.} The 368-day overlap on which \texttt{str}--\texttt{gs}
and \texttt{str}--\texttt{era5} are computed contains both a period of total
failure and a period of inverted bias. Their correlations near $r = 0.74$ are
therefore not the precision of a working instrument with an offset; they are an
average over an instrument that was working, an instrument that was dead, and an
instrument that was wrong in the opposite direction.}

\textbf{The two proxies track each other and the on-structure channel tracks
neither.} In the enlarged panel the ground station and ERA5 traces are nearly
superimposed, including on the broken days of 30 March and 6 April, while the
on-structure trace departs from both in amplitude on every day. That is the
triangulation Section~\ref{sec:absolute} quantifies, visible before any statistic
is taken, and the conclusion it supports --- that this channel is not usable as a
physical measurement without recalibration --- is reached here on the raw values
alone.

\subsection{The clock, settled two ways}

Two independent tests are available and they do not have to agree.

The \textbf{absolute} test compares each source's solar-radiation phase against
solar noon computed from the site's longitude and the equation of time. It needs
no reference sensor. The \textbf{relative} test cross-correlates each source
against the on-structure channel; it needs a reference but works on any quantity.

\begin{center}
\small
\begin{tabular}{llR{1.9cm}R{1.9cm}R{1.9cm}}
\toprule
\textbf{Source} & \textbf{Season} & \textbf{offset from solar noon} &
\textbf{shift on \texttt{tair}} & \textbf{shift on \texttt{sr}} \\
\midrule
on-structure   & winter & $+0.77$ h & --- & --- \\
on-structure   & summer & $+1.26$ h & --- & --- \\
ground station & winter & $+0.09$ h & 0 h & 1 h \\
ground station & summer & $+0.06$ h & 1 h & 1 h \\
ERA5           & winter & $+0.51$ h & 0 h & 0 h \\
ERA5           & summer & $+0.42$ h & 1 h & 1 h \\
\bottomrule
\end{tabular}
\end{center}

\panel{The on-structure radiation channel fails its own consistency check}{%
It places solar noon \emph{after} its own daily temperature maximum, which is
impossible --- the wall cannot warm before the sun reaches its highest point. Its
winter estimate has an interquartile range of \textbf{5.3 hours}, and it reads
roughly \textbf{60 W/m\textsuperscript{2} at midnight}.

\medskip
The alignment is therefore taken from \textbf{air temperature}, which all three
sources carry with good coverage and which is the quantity the rest of the study
uses. On temperature the answer is unambiguous: the ground station and ERA5 agree
with each other \emph{exactly} in both seasons, and both sit one hour from the
on-structure logger \textbf{in summer only}.

\medskip
That is the signature of a daylight-saving mismatch --- one clock advances in
spring and the other does not --- so the correction is applied by season. A
single annual figure would have been wrong in both halves of the year. After
correction the residual offset is zero on every source and season, which the
notebook asserts.}

\begin{figure}[H]
\includegraphics{P_F01_clock}
\caption{Mean diurnal solar radiation by source and season after alignment. The
on-structure channel's non-zero night-time floor is visible directly.}
\end{figure}

% ===========================================================================
\section{Absolute agreement}
\label{sec:absolute}

Three complementary views, because they answer different questions. Bias and
RMSE say how large the disagreement is. The regression slope and intercept say
what \emph{kind} it is --- a slope near one with a non-zero intercept is an
offset, a slope away from one is a scale error, and they call for different
remedies. The limits of agreement say how far a single substituted value might
fall from the one it replaces.

\begin{center}
\small
\begin{tabular}{llR{1.3cm}R{1.2cm}R{1.2cm}R{1.2cm}R{1.3cm}R{2.2cm}}
\toprule
& \textbf{pair} & \textbf{days} & \textbf{bias} & \textbf{MAE} &
\textbf{slope} & \textbf{$r$} & \textbf{limits of agr.} \\
\midrule
\multicolumn{8}{l}{\itshape air temperature [\degC]} \\
& str--gs   & 1958 & $+0.37$ & \textbf{1.12} & 0.942 & \textbf{0.988} & $-2.3$ to $+3.0$ \\
& str--era5 & 2048 & $-0.57$ & 1.53 & 0.922 & 0.978 & $-4.2$ to $+3.0$ \\
& gs--era5  & 2772 & $-0.92$ & 1.36 & 0.969 & 0.981 & $-4.0$ to $+2.1$ \\
\addlinespace
\multicolumn{8}{l}{\itshape relative humidity [\%]} \\
& str--gs   & 1926 & $-6.51$ & 8.57 & 0.775 & 0.880 & $-26$ to $+13$ \\
& str--era5 & 2048 & $-1.28$ & 9.41 & 0.686 & 0.822 & $-25$ to $+22$ \\
& gs--era5  & 2690 & $+5.80$ & 8.80 & 0.745 & 0.806 & $-17$ to $+28$ \\
\addlinespace
\multicolumn{8}{l}{\itshape solar radiation [W/m\textsuperscript{2}]} \\
& str--gs   &  368 & $-67.0$ & 153.9 & 0.577 & 0.744 & $-505$ to $+370$ \\
& str--era5 &  368 & $-43.1$ & 153.1 & 0.606 & 0.738 & $-489$ to $+403$ \\
& \textbf{gs--era5} & 2774 & $+29.1$ & \textbf{47.4} &
\textbf{0.998} & \textbf{0.930} & $-152$ to $+211$ \\
\addlinespace
\multicolumn{8}{l}{\itshape solar radiation, daylight hours only} \\
& str--gs   &  194 & $-48.8$ & 211.5 & 0.489 & 0.705 & $-600$ to $+503$ \\
& str--era5 &  195 & $-7.3$  & 210.9 & 0.494 & 0.700 & $-564$ to $+549$ \\
& gs--era5  & 1395 & $+54.1$ & \textbf{90.0} & 0.890 & \textbf{0.886} & $-190$ to $+298$ \\
\bottomrule
\end{tabular}
\end{center}

\panel{Temperature substitutes well; humidity does not}{%
\textbf{Air temperature is the one channel a proxy can stand in for.} The ground
station tracks the on-structure sensor to a mean absolute \textbf{1.12\,\degC} at
$r = 0.988$, and ERA5 to 1.53\,\degC\ at $r = 0.978$. The ground station is the
closer of the two on every measure, as its position should imply.

\medskip
\textbf{Humidity is not interchangeable.} Correlations of 0.81--0.88 and limits
of agreement spanning \textbf{fifty percentage points} mean a substituted
humidity value carries no useful precision. Any use of humidity as a regressor
must come from the source that will actually be available, not be swapped
between them.}

\panel{The radiation defect, confirmed by triangulation}{%
The on-structure radiation channel agrees with neither proxy: $r = 0.744$ against
the ground station, $0.738$ against ERA5, with mean absolute errors above
150\,W/m\textsuperscript{2} and limits of agreement spanning nearly nine hundred.

\medskip
The two proxies, meanwhile, agree with \emph{each other} at $r = 0.930$, a slope
of \textbf{0.998} and a mean absolute error of 47\,W/m\textsuperscript{2} --- a
third of the disagreement either has with the on-structure sensor, over seven
times the overlap.

\medskip
\textbf{Two independent instruments agreeing with each other and disagreeing with
a third identify which one is wrong.} Combined with the night-time floor and the
impossible solar-noon phase of Section~\ref{sec:clock}, the on-structure
radiation channel should not be used as a physical measurement without
recalibration.}

\panel{Half of a radiation record is night, and it flatters every statistic}{%
Roughly half of every solar-radiation series is a structural night-time zero.
Those zeros are real measurements and are removed nowhere in this study, but a
statistic computed over the whole day is computed largely on a block of samples
on which all three sources agree trivially, because all three are reporting the
same zero.

\medskip
Restricting to the hours the sun stands above the horizon --- a geometric mask
taken from the site coordinates, independent of all three sources --- removes
that flattery, and the effect is large. The correlation between the two proxies
falls from \textbf{0.930 to 0.886} and their mean absolute difference nearly
doubles, from \textbf{47.4 to 90.0}\,W/m\textsuperscript{2}; their regression
slope falls from 0.998 to 0.890. The two comparisons involving the on-structure
channel move less, from $r \approx 0.74$ to $r \approx 0.70$, because that
channel disagrees with both proxies during the day as well as at night.

\medskip
\textbf{The daylight figures are the ones that describe how well a proxy could
stand in for a radiation measurement}, since a reconstruction is only ever
tested where there is something to reconstruct. The all-hours figures are kept
beside them because they are what a naive comparison reports.

\medskip
The mask is geometry, not radiation. It takes no account of cloud, and none of
the local horizon formed by the terrain around the site, so it marks when the
sun is up rather than when the wall was lit.}

\begin{figure}[H]
\includegraphics{P_F02_agreement_tair}
\caption{Air temperature. Top: each proxy against the on-structure sensor, with
the identity line. Bottom: Bland--Altman, difference against mean, with the mean
difference and the limits of agreement.}
\end{figure}

% ===========================================================================
\section{Dynamic agreement}
\label{sec:dynamic}

A proxy can sit two degrees low and still drive the wall correctly, or match on
average and miss every daily swing.

\subsection{The daily cycle}

\begin{center}
\small
\begin{tabular}{llR{2.1cm}R{2.1cm}R{1.5cm}}
\toprule
& \textbf{Source} & \textbf{diurnal amplitude} & \textbf{ratio to
on-structure} & \textbf{peak hour} \\
\midrule
\multicolumn{5}{l}{\itshape air temperature [\degC]} \\
& on-structure   & 3.31 & 1.00 & 13 \\
& ground station & 3.80 & \textbf{1.15} & 13 \\
& ERA5           & 3.94 & \textbf{1.19} & 14 \\
\addlinespace
\multicolumn{5}{l}{\itshape relative humidity [\%]} \\
& on-structure   & 10.22 & 1.00 &  7 \\
& ground station & 11.87 & 1.16 &  6 \\
& ERA5           & 13.91 & 1.36 &  5 \\
\addlinespace
\multicolumn{5}{l}{\itshape solar radiation [W/m\textsuperscript{2}], daylight hours} \\
& on-structure   & 291.9 & 1.00 & 12 \\
& ground station & 243.0 & 0.83 & 12 \\
& ERA5           & 282.6 & 0.97 & 12 \\
\bottomrule
\end{tabular}
\end{center}

The radiation rows are computed over the daylight hours and carry no trough
hour. Over the full day the trough of a radiation cycle is an arbitrary point
inside the night --- the all-hours computation returns hours 23, 1 and 2 for the
three sources, which are three different ways of saying ``dark'' --- and a
number that arbitrary does not belong in a comparison. All three sources peak at
hour 12, which is the one phase statement the radiation channels agree on.

Radiation is also the one quantity where the on-structure sensor shows the
\emph{largest} daily amplitude rather than the smallest. That is consistent with
the defect established above rather than with the damping seen in temperature: a
channel reading roughly 60\,W/m\textsuperscript{2} at midnight and disagreeing
with both proxies has no claim to a physically meaningful amplitude.

\panel{The expectation was that ERA5 would flatten the daily cycle. It does the
opposite}{%
A nine-kilometre grid average should smooth a microclimate away, so the
reanalysis was expected to show the \emph{smallest} diurnal amplitude. It shows
the \textbf{largest} --- 19\,\% above the on-structure sensor, with the ground
station 15\,\% above.

\medskip
The sensor closest to the wall is the most \textbf{damped} of the three, and on
reflection that is what the physics requires. The on-structure channel measures
air inside an instrument housing mounted on a massive masonry wall. The housing
has thermal mass and the wall behind it has a great deal more, so the air the
sensor sees is buffered against the daily swing that a freely ventilated screen
in a field records in full.

\medskip
\textbf{The consequence for imputation is favourable and was not anticipated.} A
proxy that under-represented the daily cycle would have been unable to
reconstruct it. Both proxies over-represent it instead, which is a scale
difference --- correctable by regression --- rather than missing information.}

ERA5's temperature peak falls one hour later than the other two. The humidity
peaks spread over three hours, which is a further reason not to substitute
humidity between sources.

\begin{figure}[H]
\includegraphics{P_F03_diurnal_tair}
\caption{Mean diurnal cycle of air temperature by source. Left: absolute. Right:
departure from each source's own daily mean, which isolates the shape.}
\end{figure}

\subsection{What survives differencing}

The wall responds to change, not to level, and the first difference is the
quantity invariant to every arbitrary constant in the chain.

\begin{center}
\small
\begin{tabular}{llR{1.6cm}R{1.8cm}R{1.5cm}R{1.4cm}}
\toprule
& \textbf{pair} & \textbf{$r$ levels} & \textbf{$r$ differences} &
\textbf{lost} & \textbf{slope} \\
\midrule
\multicolumn{6}{l}{\itshape air temperature} \\
& str--gs   & 0.988 & \textbf{0.699} & $-0.29$ & 0.684 \\
& str--era5 & 0.978 & \textbf{0.690} & $-0.29$ & 0.639 \\
& gs--era5  & 0.981 & 0.731 & $-0.25$ & 0.679 \\
\addlinespace
\multicolumn{6}{l}{\itshape relative humidity} \\
& str--gs   & 0.880 & 0.546 & $-0.33$ & 0.542 \\
& str--era5 & 0.822 & 0.489 & $-0.33$ & 0.498 \\
\addlinespace
\multicolumn{6}{l}{\itshape solar radiation} \\
& str--gs   & 0.744 & 0.491 & $-0.25$ & 0.343 \\
& gs--era5  & 0.930 & 0.648 & $-0.28$ & 0.537 \\
\bottomrule
\end{tabular}
\end{center}

\panel{Most of the level agreement is a shared annual cycle}{%
Differencing costs about \textbf{0.29} of correlation on every temperature pair
and about 0.33 on humidity. Two sources correlating at 0.99 in levels agree only
at 0.70 on what happened between one hour and the next.

\medskip
This is not a defect in either instrument. Two thermometers a few kilometres
apart share the season and the synoptic weather almost perfectly, and differ on
the local gust, the passing cloud and the sensor's own response time --- all of
which live in the difference.

\medskip
\textbf{It sets the ceiling for proxy-driven reconstruction.} A model driven by a
proxy can recover the part of the wall's movement that the proxy shares with the
on-structure sensor, and that share is far smaller in the domain the wall
actually responds to. The slopes below one --- 0.64 to 0.68 for temperature ---
are partly regression dilution, since the reference carries measurement noise of
its own, so they bound rather than measure the attenuation.}

% ===========================================================================\section{Relationship with the inclination}
\label{sec:target}

Every statistic up to this point was computed on the joined frame: 26 July 2018
to 13 August 2026, both instrument eras, and ragged throughout. No correlation
against the target is computed that way below, for two reasons that this section
establishes before it reports anything.

\subsection{The windows this comparison can use}
\label{sec:tiers}

The prediction studies never correlate against the target on a span of that
kind. They cut contiguous blocks of near-complete coverage, group them into
tiers according to which channels a tier requires, and report every correlation,
slope and operator inside a tier. This report now does the same, with one
addition the sibling studies do not need: because it spans both eras, it must
also keep them apart. The instrument at \texttt{st02} was physically reinstalled
on 21 February 2025 and the two packages share no baseline, so no window here
crosses that date.

\begin{center}
\small
\begin{tabularx}{\textwidth}{l L{3.4cm} R{1.4cm} R{2.6cm} R{1.9cm}}
\toprule
\textbf{Tier} & \textbf{Requires} & \textbf{Blocks} & \textbf{Longest block} &
\textbf{Total days} \\
\midrule
P0 & target and all three sources & 7 & 86.6 d & 343.4 \\
P1 & target and both proxies     & 7 & 86.6 d & 386.4 \\
P2 & target and both proxies, legacy era & 23 & 107.8 d & 1271.2 \\
\bottomrule
\end{tabularx}
\end{center}

\noindent{\footnotesize Blocks are contiguous runs in which every required
channel is present, interruptions up to six hours absorbed, minimum length
twenty days --- the same parameters the prediction studies use, so the
inventories are comparable. Full listing in \texttt{P\_13}.}

\medskip
\textbf{The radiation record sets the shape of the tiers by itself.}
\texttt{sr\_str} begins on 21 February 2025, which is the era boundary, so a
tier requiring the on-structure pyranometer cannot reach the legacy era at all.
The three-source radiation comparison therefore lives entirely inside the
current era, while the two proxies additionally carry a legacy era the
on-structure channel can never enter. That is the whole of what ``solar
radiation is available for only a fraction of the period'' means here, and it is
a property of the tiers rather than a caveat attached to them.

Tiers P0 and P1 differ only in whether the on-structure channels are required.
P1 contains every hour of P0 and a little more, so the two are not independent
samples and their agreement is weaker evidence than it looks. Where they differ
at all below, the difference is one block.

\subsection{A defect that had to be settled first}
\label{sec:faults}

Blocking alone did not reconcile this report with the prediction studies, and
the reason turned out to be in the record rather than in the method.

\panel{Single-hour acquisition faults, and why they had to be registered}{%
The current era contains hours in which the calibrated inclination moves by
hundreds of millidegrees and returns on the following sample. At
\textbf{2025-05-10 14:00} the raw inclination moves $-220$\,mdeg, the air
temperature rises $5.9\,\degC$ and the battery voltage dips --- all within the
same hour, all recovering at the next. Three physically independent channels do
not fail together for one sample because a wall moved.

\medskip
\textbf{The compensation is not the cause and differencing is not the cause.}
The raw channel already carries 110 such hours in the current era against 47 in
the legacy era; the calibrated channel carries 166 against 37, inheriting them
and enlarging those where the temperature spiked as well. The index is a
complete hourly grid, so no difference is ever taken across a gap.

\medskip
They matter because they live almost entirely in the first difference: one
unreversed hour contributes two large differences of opposite sign. Left in
place they do not merely add noise --- they reverse conclusions.

\medskip
\textbf{The hours are marked, not deleted.} A rolling median and a median
absolute deviation over a one-day half-window score every hour, and those beyond
eight robust standard deviations are entered in a register,
\texttt{P\_23}, together with the evidence for each: the signed step, how many
other channels were simultaneously anomalous, and whether the excursion reversed
on the next hour. The register holds \textbf{175 hours, 0.248\,\%} of the grid.
The mask is applied to the \emph{level}, so one decision governs the level
statistics and the differenced ones rather than two separate rules, and a
sustained shift --- a real re-levelling --- is not caught by it.}

\begin{center}
\small
\begin{tabular}{lR{2.0cm}R{2.0cm}R{2.4cm}}
\toprule
& \textbf{$r$ unscreened} & \textbf{$r$ screened} & \textbf{sibling study} \\
\midrule
\multicolumn{4}{l}{\itshape tier P0, on-structure air temperature} \\
\quad levels      & $-0.757$ & $\mathbf{-0.952}$ & $-0.956$ \\
\quad differences & $-0.148$ & $\mathbf{-0.840}$ & $-0.886$ \\
\addlinespace
\multicolumn{4}{l}{\itshape tier P0, on-structure solar radiation} \\
\quad differences & $\mathbf{+0.036}$ & $\mathbf{-0.533}$ & $-0.624$ \\
\addlinespace
\multicolumn{4}{l}{\itshape tier P2, on-structure air temperature} \\
\quad differences & $-0.890$ & $-0.890$ & $-0.901$ \\
\bottomrule
\end{tabular}
\end{center}

\noindent{\footnotesize Sibling figures from
\texttt{studies/inclination\_prediction} for the current era and
\texttt{studies/inclination\_prediction\_legacy} for the legacy era.}

\medskip
\textbf{The register changes the current era and leaves the legacy era alone.}
Every legacy figure moves by 0.002 or less, which is why tier P2 reproduced its
sibling study before any screening at all; the legacy faults exist but are fifty
times smaller, a maximum step of 51\,mdeg against 1314. In the current era the
correction is decisive in both domains, and it recovers the sibling studies'
values from a record that had appeared to contradict them.

\panel{A reversal that was not a defect}{%
Unscreened, the on-structure radiation channel correlates \textbf{positively}
with the target in first differences, $r = +0.036$, while both proxies stay
negative. Read on its own that is a fourth symptom of a channel already known to
be defective.

\medskip
It is not. Screened, the same channel gives $r = -0.533$, in agreement with both
proxies and with the sibling study. \textbf{The apparent sign reversal was
produced entirely by the acquisition faults}, and nothing in this report reopens
the sign question of \texttt{docs/raw-data-format.md} §7.5.}

\subsection{Correlation structure}
\label{sec:corr}

With the windows defined and the faults registered, the correlation matrix over
the target and every channel a tier carries, on levels and on first differences.

\begin{center}
\small
\begin{tabular}{llR{1.9cm}R{1.9cm}R{1.9cm}}
\toprule
& & \textbf{on-structure} & \textbf{ground station} & \textbf{ERA5} \\
\midrule
\multicolumn{5}{l}{\itshape air temperature} \\
& P0 levels      & $-0.952$ & $-0.940$ & $-0.932$ \\
& P0 differences & $-0.840$ & $-0.582$ & $-0.570$ \\
& P2 levels      & $-0.541$ & $-0.543$ & $-0.539$ \\
& P2 differences & $-0.890$ & $-0.668$ & $-0.569$ \\
\addlinespace
\multicolumn{5}{l}{\itshape solar radiation} \\
& P0 levels      & $-0.142$ & $-0.442$ & $-0.460$ \\
& P0 differences & $-0.533$ & $-0.430$ & $-0.527$ \\
& P2 levels      & ---      & $-0.222$ & $-0.215$ \\
& P2 differences & ---      & $-0.491$ & $-0.580$ \\
\bottomrule
\end{tabular}
\end{center}

Correlations between two trending series are inflated by the shared trend and
say little about a driver--response relationship; the differenced matrix is the
honest one for a signal with drift. Both are shown so the gap is visible, which
is the convention the prediction studies use.

\textbf{The two eras answer differently, and the difference is instructive.} In
the current era the level correlation is the stronger one for temperature and
the differenced correlation the weaker; in the legacy era it is the reverse. The
legacy blocks are drawn from seven calendar years and no single block spans an
annual cycle, so their levels carry less common trend to inflate them, while
their hour-to-hour movement is measured on a quieter instrument.

\textbf{Radiation is the one quantity that gains under differencing in both
eras.} Its level correlation with the target is weak everywhere --- between
$-0.14$ and $-0.46$ --- and its differenced correlation is stronger for every
source in every tier. Radiation drives the daily cycle and carries almost none
of the seasonal level, which is exactly the signature differencing exposes.

\begin{figure}[H]
\includegraphics{P_F06_correlations}
\caption{Correlation structure by tier, on levels (left) and on first
differences (right). The block of high mutual correlation among the drivers
themselves is quantified in Section~\ref{sec:contribution}.}
\end{figure}

\subsection{Relationship with the target, by tier}
\label{sec:byTier}

Each source against the one target, in both domains, inside each tier's blocks.
The slope describes the residual response that survives the calibration;
\textbf{it is not a candidate compensation coefficient}.

\begin{center}
\small
\begin{tabular}{llR{1.5cm}R{2.3cm}R{1.3cm}R{1.2cm}}
\toprule
& \textbf{Source} & \textbf{$r$} & \textbf{slope
[\,mdeg\,unit$^{-1}$]} & \textbf{s.e.} & \textbf{$R^2$} \\
\midrule
\multicolumn{6}{l}{\itshape air temperature [\,mdeg\,\degC$^{-1}$], tier P0, levels} \\
& on-structure   & $-0.952$ & $-4.602$ & 0.067 & 0.907 \\
& ground station & $-0.940$ & $-4.842$ & 0.074 & 0.883 \\
& ERA5           & $-0.932$ & $-4.766$ & 0.074 & 0.868 \\
\addlinespace
\multicolumn{6}{l}{\itshape air temperature, tier P0, first differences} \\
& on-structure   & $-0.840$ & $-3.010$ & 0.041 & 0.706 \\
& ground station & $-0.582$ & $-2.391$ & 0.050 & 0.339 \\
& ERA5           & $-0.570$ & $-2.336$ & 0.051 & 0.325 \\
\addlinespace
\multicolumn{6}{l}{\itshape solar radiation [\,mdeg\,(W\,m$^{-2}$)$^{-1}$], tier P0, first differences} \\
& on-structure   & $-0.533$ & $-0.0167$ & 0.0006 & 0.284 \\
& ground station & $-0.430$ & $-0.0187$ & 0.0008 & 0.185 \\
& ERA5           & $-0.527$ & $-0.0276$ & 0.0008 & 0.278 \\
\addlinespace
\multicolumn{6}{l}{\itshape air temperature, tier P2, first differences} \\
& on-structure   & $-0.890$ & $-2.653$ & 0.020 & 0.792 \\
& ground station & $-0.668$ & $-1.966$ & 0.024 & 0.447 \\
& ERA5           & $-0.569$ & $-1.766$ & 0.023 & 0.323 \\
\bottomrule
\end{tabular}
\end{center}

\noindent{\footnotesize Standard errors are Newey--West at 24 lags. The full
table, every quantity, every tier, both domains, and both the screened and the
unscreened frame, is \texttt{P\_16}.}

\panel{Every source agrees on the sign, and the proxies pay a measurable price}{%
All twelve correlations above are \textbf{negative}, which is the direction the
site's geometry requires: the valley face is the sun-exposed one, it expands
more under heating, and the wall tips toward the mountain
(\texttt{docs/raw-data-format.md} §7.5). No source and no tier reverses it.

\medskip
\textbf{The on-structure sensor retains substantially more of the differenced
relationship than either proxy} --- $-0.840$ against $-0.582$ and $-0.570$ in
the current era, $-0.890$ against $-0.668$ and $-0.569$ in the legacy era. That
gap is the price of substitution, and it is the quantity a reconstruction study
inherits: a proxy-driven model can recover the part of the wall's movement that
the proxy shares with the wall's own sensor, and in the difference domain that
share is roughly two thirds.

\medskip
\textbf{Solar radiation reaches roughly half the differenced correlation of air
temperature}, $-0.53$ against $-0.84$ in tier P0, from sources that agree with
each other far better than the radiation channels do. It is a real driver in the
domain the wall responds to, not a residual.}

\begin{figure}[H]
\includegraphics{P_F07_relationship_tiers}
\caption{Slope against the calibrated inclination, by tier and domain, for air
temperature and solar radiation. Error bars are Newey--West standard errors at
24 lags.}
\end{figure}

\subsection{Transport delay and thermal inertia, on two bands}
\label{sec:operators}

Before a driver is compared it is given the chance to act through the operator
that suits it. A \textbf{transport delay} shifts the driver without changing its
shape; a \textbf{thermal inertia} low-passes it through a single-pole filter.
The two are scanned jointly over 13 delays and 16 time constants, because
scanning either alone can attribute to one what belongs to the other, and the
instantaneous case stays in the grid so that it competes rather than being
assumed away.

The scan runs on each tier's \emph{longest single block}, not on the union of
its blocks. A delay is applied with a shift, a time constant with a recursive
filter and the band limit with a rolling mean; all three read across adjacent
rows, so a union would let a filter draw values from the far side of a
months-long interruption.

\begin{center}
\small
\begin{tabular}{llR{1.3cm}R{1.2cm}R{1.5cm}R{1.4cm}R{2.0cm}}
\toprule
\textbf{Tier} & \textbf{Driver} & \textbf{delay} & \textbf{$\tau$} &
\textbf{$r$} & \textbf{$R^2$} & \textbf{$R^2$ at $d=\tau=0$} \\
\midrule
\multicolumn{7}{l}{\itshape diurnal band --- carried forward} \\
P0 & \texttt{tair\_str}  & 0 h & 0 h & $-0.977$ & 0.955 & 0.955 \\
P0 & \texttt{tair\_gs}   & 0 h & 0 h & $-0.930$ & 0.865 & 0.865 \\
P0 & \texttt{tair\_era5} & 0 h & 0 h & $-0.916$ & 0.839 & 0.839 \\
P0 & \texttt{sr\_str}    & 0 h & 1 h & $-0.542$ & 0.294 & 0.262 \\
P0 & \texttt{sr\_gs}     & 0 h & 2 h & $-0.841$ & \textbf{0.708} & 0.570 \\
P0 & \texttt{sr\_era5}   & 0 h & 1 h & $-0.829$ & 0.686 & 0.618 \\
\addlinespace
P2 & \texttt{tair\_str}  & 0 h & 0 h & $-0.964$ & 0.930 & 0.930 \\
P2 & \texttt{sr\_gs}     & 0 h & 2 h & $-0.746$ & 0.556 & 0.437 \\
P2 & \texttt{sr\_era5}   & 0 h & 1 h & $-0.711$ & 0.506 & 0.460 \\
\bottomrule
\end{tabular}
\end{center}

\noindent{\footnotesize Delays bounded at twelve hours per
\texttt{docs/raw-data-format.md} §7.5: both air temperature and solar radiation
are external forcings, and beyond half a diurnal cycle a delay is
indistinguishable from a lead. Tier P1 reproduces P0 exactly, because the two
share their longest block. Full-band results and the complete table are
\texttt{P\_17} and \texttt{P\_18}.}

\panel{No driver needs a transport delay; only radiation needs an inertia}{%
\textbf{Every optimum sits at zero delay}, on both bands and in both eras. A
grid cell averaged over nine kilometres was expected to show a different
apparent lag from a sensor on the wall; it does not. A proxy therefore needs no
phase correction before use, which is a simplification the scan earned rather
than an assumption.

\medskip
\textbf{Air temperature also needs no inertia} --- its optimum is
$\tau = 0$ for every source in every tier, and the operator adds nothing to the
instantaneous fit.

\medskip
\textbf{Solar radiation does.} Both proxies optimise at a one- or two-hour time
constant, and the gain is not marginal: for the ground station explained
variance rises from \textbf{0.570 to 0.708}. Radiation heats a surface which
then warms the air and the masonry, so a short lag between the forcing and the
response it produces is what the physics predicts. This is the one place in the
report where fitting an operator changes a conclusion.

\medskip
\textbf{The full band produced no boundary optimum.} The prediction study found
solar radiation pinned at the last time constant in the grid, 168 hours, which
it correctly diagnosed as a filter turning a daily cycle into a seasonal
envelope. On these blocks that does not occur for any driver: no full-band
optimum sits at the grid edge. The blocks here are shorter than an annual cycle,
so there is no seasonal component for a week-wide filter to reach for.}

\begin{figure}[H]
\includegraphics{P_F08_operators_P0}
\caption{Explained variance over the delay-by-inertia grid, diurnal band, tier
P0. Every optimum sits on the zero-delay edge; the radiation panels brighten one
to two hours into the inertia axis, which is the operator the table reports.}
\end{figure}

\subsection{Does any radiation channel lead the response?}
\label{sec:signed}

The scan above is bounded at non-negative delays because a forcing must precede
the response it causes. That bound is right for \emph{use} --- a lead cannot be
applied by a forecasting system, which would need the driver's future values at
the moment the forecast is issued --- but it also hides a diagnostic. If an
optimum genuinely sat at a negative delay, the bounded scan would report the
boundary instead of the anomaly. Scanning signed delays makes it visible.

\begin{center}
\small
\begin{tabular}{llR{1.9cm}R{2.0cm}R{2.0cm}}
\toprule
\textbf{Tier} & \textbf{Channel} & \textbf{peak lag} & \textbf{$r$ at peak} &
\textbf{$r$ at lag 0} \\
\midrule
P0 & \texttt{sr\_str}  & $+1$ h & $-0.524$ & $-0.512$ \\
P0 & \texttt{sr\_gs}   & $+1$ h & $-0.825$ & $-0.755$ \\
P0 & \texttt{sr\_era5} & $+1$ h & $-0.821$ & $-0.786$ \\
P2 & \texttt{sr\_gs}   & $+1$ h & $-0.721$ & $-0.661$ \\
P2 & \texttt{sr\_era5} & $+1$ h & $-0.700$ & $-0.679$ \\
\bottomrule
\end{tabular}
\end{center}

\textbf{Every radiation channel peaks at a positive lag of one hour}, in both
eras and including the defective on-structure channel. A positive lag means the
driver leads the response, which is the only direction a forcing can act in, and
one hour is consistent with the one- to two-hour inertia the operator scan
selected independently. No channel leads the deformation, so the causal
constraint imposed on every model in this project costs nothing here, and the
sign anomaly of Section~\ref{sec:faults} is confirmed as an artefact of the
acquisition faults rather than a property of any instrument.

\subsection{What radiation adds once temperature is present}
\label{sec:contribution}

Radiation heats the air, so a marginal correlation between radiation and the
target partly re-measures the temperature relationship. The variance inflation
factor quantifies the duplication, and a joint fit reports what radiation adds
beyond temperature alone.

\begin{center}
\small
\begin{tabular}{llR{1.2cm}R{1.6cm}R{1.5cm}R{1.5cm}R{1.4cm}}
\toprule
\textbf{Tier} & \textbf{Source} & \textbf{VIF} &
\textbf{$r$ between} & \textbf{$R^2$ \texttt{tair}} &
\textbf{$R^2$ joint} & \textbf{$\Delta R^2$} \\
\midrule
\multicolumn{7}{l}{\itshape first differences} \\
P0 & on-structure   & 1.37 & 0.520 & 0.706 & 0.719 & $+0.013$ \\
P0 & ground station & 1.57 & 0.602 & 0.339 & 0.349 & $+0.010$ \\
P0 & ERA5           & 2.01 & 0.709 & 0.325 & 0.355 & $+0.031$ \\
P2 & ground station & 1.64 & 0.626 & 0.447 & 0.456 & $+0.009$ \\
P2 & ERA5           & 1.67 & 0.634 & 0.323 & 0.404 & $\mathbf{+0.081}$ \\
\bottomrule
\end{tabular}
\end{center}

\noindent{\footnotesize Newey--West standard errors at 24 lags. The radiation
term clears $|t| = 6$ in every row above and reaches $|t| = 45$ in the last.
Level-domain rows are in \texttt{P\_20}.}

\panel{Radiation earns a place, and it is a small one}{%
\textbf{Collinearity is mild.} Variance inflation factors run from 1.37 to 2.01,
so the two drivers are not interchangeable and a joint fit is interpretable.

\medskip
\textbf{The increment is statistically overwhelming and practically small.}
Every radiation term clears $|t| = 6$, which on tens of thousands of hourly
samples establishes only that the effect is not zero. What it buys is one to
three points of explained variance in the current era, and eight in the legacy
era for ERA5.

\medskip
\textbf{The recommendation follows the size, not the significance.} A radiation
channel is worth carrying as a second regressor where it is free --- both
proxies supply one --- but it does not change what a temperature-driven model
can do, and it is not a substitute for temperature anywhere.}

\subsection{What the pooled span cost}
\label{sec:pooled}

The previous version of this section reported the relationship on the whole
record at once. Against the tiered, screened figures:

\begin{center}
\small
\begin{tabular}{llR{2.2cm}R{2.2cm}}
\toprule
& & \textbf{pooled span} & \textbf{tier P0} \\
\midrule
\multicolumn{4}{l}{\itshape air temperature, first differences} \\
& on-structure   & $-0.155$ & $\mathbf{-0.840}$ \\
& ground station & $-0.091$ & $-0.582$ \\
& ERA5           & $-0.084$ & $-0.570$ \\
\addlinespace
\multicolumn{4}{l}{\itshape solar radiation, first differences} \\
& on-structure   & $+0.087$ & $\mathbf{-0.533}$ \\
& ground station & $-0.088$ & $-0.430$ \\
& ERA5           & $-0.091$ & $-0.527$ \\
\bottomrule
\end{tabular}
\end{center}

\textbf{The earlier conclusion was an artefact of the analysis window and the
unregistered faults, not a finding about the wall.} That version reported the
differenced relationship as weak for every source, with $R^2$ between 0.007 and
0.024, and read it as agreement with the prediction studies that the drivers
contribute nothing measurable. The prediction studies report no such thing: on
their own blocks they find air temperature at $r \approx -0.89$ in first
differences, which is what this report now also finds. The pooled figures are
retained in \texttt{P\_08} as the contrast, and are not interpreted anywhere.

% ===========================================================================
\section{Verdict}
\label{sec:verdict}

\subsection{Availability, which decides everything else}

\begin{center}
\small
\begin{tabular}{llR{2.4cm}R{3.2cm}}
\toprule
& \textbf{Source} & \textbf{coverage of archive span} & \textbf{present when the
target is missing} \\
\midrule
\multicolumn{4}{l}{\itshape air temperature} \\
& on-structure   & 70.1\,\% & \textbf{0.0\,\%} \\
& ground station & 94.8\,\% & \textbf{94.3\,\%} \\
& ERA5           & 99.4\,\% & \textbf{99.6\,\%} \\
\addlinespace
\multicolumn{4}{l}{\itshape solar radiation} \\
& on-structure   & 12.7\,\% & \textbf{0.0\,\%} \\
& ground station & 94.9\,\% & \textbf{94.4\,\%} \\
& ERA5           & 99.4\,\% & \textbf{99.6\,\%} \\
\bottomrule
\end{tabular}
\end{center}

\panel{The deadlock is broken}{%
The on-structure sensor is present on \textbf{none} of the hours the target is
missing --- that is what outage-driven missingness means, and it is why the
previous study had to refuse the driver-based filler.

\medskip
The ground station is present on \textbf{94.3\,\%} of those hours and ERA5 on
\textbf{99.6\,\%}. The availability criterion that emptied the candidate pool is
satisfied by both proxies, and comfortably.

\medskip
\textbf{Radiation behaves the same way, and the asymmetry is extreme.} The
on-structure pyranometer covers \textbf{12.7\,\%} of the archive span and, like
every on-structure channel, none of the hours the target is missing. Both
proxies supply radiation on more than 94\,\% of exactly those hours. A radiation
regressor is therefore available for reconstruction only if it comes from a
proxy, which is the same conclusion temperature reached and for the same
reason.}

\subsection{Which source, for what}

\begin{center}
\small
\begin{tabularx}{\textwidth}{L{3.2cm}Y}
\toprule
\textbf{Use} & \textbf{Recommendation} \\
\midrule
Gap reconstruction & \textbf{ERA5} as primary, at 99.6\,\% availability, with the
ground station as a cross-check where both exist. ERA5's slightly worse agreement
(1.53 against 1.12\,\degC) is outweighed by its near-total coverage. \\
Regressor for a forecast model & \textbf{Either}, with ERA5 preferred for
continuity. Both need a scale correction: their diurnal amplitude runs 15--19\,\%
above the wall's own sensor. \\
Compensation & \textbf{On-structure only}, always. The calibration is an
instrument correction and must use the temperature at the transducer. \\
Humidity, any use & \textbf{Do not substitute.} Limits of agreement span fifty
percentage points. \\
Radiation as a second regressor & \textbf{Either proxy}, applied through a
one- to two-hour inertia filter. It adds one to three points of explained
variance over air temperature alone, which is worth taking where it is free and
is not worth designing around. \\
Radiation, as a measurement & \textbf{Not the on-structure channel} until it is
recalibrated. The two proxies agree with each other and both disagree with it,
and it reaches barely half their differenced correlation with the target. \\
\bottomrule
\end{tabularx}
\end{center}

\subsection{The feature specification for the imputation study}

\begin{center}
\small
\begin{tabular}{llR{1.1cm}R{1.0cm}R{1.5cm}R{1.6cm}R{1.9cm}}
\toprule
\textbf{Channel} & \textbf{Source} & \textbf{delay} & \textbf{$\tau$} &
\textbf{$R^2$ diurnal} & \textbf{$r$ diff.} & \textbf{available in gaps} \\
\midrule
\texttt{tair\_str}  & on-structure   & 0 h & 0 h & 0.955 & $-0.840$ & 0.0\,\% \fail \\
\texttt{tair\_gs}   & ground station & 0 h & 0 h & 0.865 & $-0.582$ & 94.3\,\% \pass \\
\texttt{tair\_era5} & ERA5           & 0 h & 0 h & 0.839 & $-0.570$ & 99.6\,\% \pass \\
\addlinespace
\texttt{sr\_str}    & on-structure   & 0 h & 1 h & 0.294 & $-0.533$ & 0.0\,\% \fail \\
\texttt{sr\_gs}     & ground station & 0 h & 2 h & 0.708 & $-0.430$ & 94.4\,\% \pass \\
\texttt{sr\_era5}   & ERA5           & 0 h & 1 h & 0.686 & $-0.527$ & 99.6\,\% \pass \\
\bottomrule
\end{tabular}
\end{center}

\noindent{\footnotesize Operators from the diurnal-band scan on tier P0's
longest block; differenced correlations from tier P0. Full table in
\texttt{P\_12}.}

\medskip
Both proxies enter the imputation study \textbf{unlagged}, which is a
simplification the operator scan earned rather than an assumption. Their
temperature channels enter without an inertia filter as well; their
\textbf{radiation channels do not} --- a one- to two-hour time constant raises
explained variance for the ground station from 0.570 to 0.708 and is the one
operator in this report that changes a conclusion.

% ===========================================================================
\section{Caveats}

\textbf{Agreement is not interchangeability.} Two temperature series can agree to
1.1\,\degC\ and still imply different things for a wall, because the wall
responds to the difference and the difference agrees far less well.

\textbf{The differenced slopes are attenuated.} Regression dilution biases them
toward zero whenever the reference carries noise of its own, which it does. They
bound the attenuation rather than measuring it.

\textbf{The radiation comparison rests on a short and partly broken overlap.}
On-structure radiation exists only from 2025-02-21, giving 368 days against the
2774 available for the two proxies, and 194 days once the daylight mask is
applied. Section~\ref{sec:series} shows that the channel failed outright for two
months of that overlap and under-read for four more, so every statistic
involving it pools three different instrument states and should be read as a
description of the record rather than of a sensor.

\textbf{The legacy tier measures a different instrument.} Tier P2 lies entirely
before the 21 February 2025 reinstallation. The package installed on that date
shares no baseline with the one it replaced, so a result obtained on P2
describes the earlier instrument. No window in Section~\ref{sec:target} spans
the boundary, and no figure is carried across it.

\textbf{Tiers P0 and P1 overlap.} P1 contains every hour of P0 and one further
block. Agreement between the two is therefore weaker evidence than agreement
between independent samples, and the two differ nowhere that matters in this
report because they share their longest block.

\textbf{The daylight mask is geometry, not radiation.} It marks the hours the
sun stands above the horizon at the site's coordinates. It takes no account of
cloud, and none of the local horizon formed by the terrain around the site, so
it says when the sun was up rather than when the wall was lit.

\textbf{The fault register is a threshold, and thresholds can be argued with.}
Eight robust standard deviations is deliberately conservative; the count of
registered hours changes with it, though the conclusions do not over the range
tested. The register is published as \texttt{P\_23} with the evidence for every
flagged hour precisely so that the judgement can be inspected rather than taken
on trust. The faults themselves are a property of
\texttt{data/interim/unified/}, which every study in this project reads, and
correcting them at that source rather than in this study would be the more
durable fix.

\textbf{One station.} All three sources are compared at \texttt{st02} only, by
instruction and for comparability with the prediction studies.

\textbf{The clock correction is empirical.} It is measured on temperature and
applied by season. It reconciles the sources to within one hour, which is the
grid resolution; a sub-hourly residual would not be visible.

% ===========================================================================
\section{Artefact inventory}

\begin{center}
\small
\begin{tabularx}{\textwidth}{L{3.0cm}Y}
\toprule
\textbf{Artefact} & \textbf{Content} \\
\midrule
\texttt{P\_01} -- \texttt{P\_03} & Channel inventory, clock test, comparable
channels with overlaps \\
\texttt{P\_04}, \texttt{P\_05} & Pairwise agreement; the same split by season and
hour \\
\texttt{P\_06}, \texttt{P\_07} & Diurnal amplitude and phase; first-difference
agreement \\
\texttt{P\_08}, \texttt{P\_10} & Relationship with the target on the pooled
span, retained as the contrast case; pooled slopes \\
\texttt{P\_11}, \texttt{P\_12} & Verdicts; the feature specification \\
\texttt{P\_13} & Block inventory: the windows each tier can use \\
\texttt{P\_14}, \texttt{P\_15} & Correlation structure per tier, levels and
first differences \\
\texttt{P\_16} & Relationship with the target by tier and domain, screened and
unscreened \\
\texttt{P\_17}, \texttt{P\_18} & Operator scans, full band and diurnal band \\
\texttt{P\_19} & Signed-delay control on the radiation channels \\
\texttt{P\_20} & What radiation adds once temperature is present \\
\texttt{P\_21} & Radiation agreement over the daylight hours \\
\texttt{P\_22} & Radiation and temperature availability, including during
target gaps \\
\texttt{P\_23} & The acquisition-fault register, with the evidence per hour \\
\addlinespace
\texttt{P\_F09} & Solar radiation as measured, per source, over the
on-structure record, with a fortnight enlarged \\
\texttt{P\_F01} & Diurnal radiation by source and season, after alignment \\
\texttt{P\_F02} & Scatter and Bland--Altman, air temperature \\
\texttt{P\_F03} & Mean diurnal cycle by source \\
\texttt{P\_F05} & Correlation and slope against the target, pooled span \\
\texttt{P\_F06} & Correlation heatmaps by tier \\
\texttt{P\_F07} & Slope against the target by tier and domain \\
\texttt{P\_F08} & Operator scans, diurnal band, one figure per tier \\
\bottomrule
\end{tabularx}
\end{center}

The canonical channel reference is \texttt{docs/proxy-data-dictionary.md},
generated from the mapping tables in \texttt{pc\_lib} so that code and
documentation cannot drift apart.

Reproduction, from the study directory with \texttt{neuralprophet\_env} active:

\begin{quote}
\ttfamily\small
jupytext -{}-to ipynb proxy\_comparison\_study.py \\
jupyter nbconvert -{}-to notebook -{}-execute -{}-inplace \\
\phantom{jupyter nbconvert }proxy\_comparison\_study.ipynb
\end{quote}

\texttt{studies/unified\_dataset/build\_unified\_dataset.py} must have run first;
this study reads its output and the two proxy exports, and nothing else.

\end{document}
